%%%% LISTA DE ABREVIATURAS E SIGLAS 
%%
%% Relação, em ordem alfabética, das abreviaturas (representação de uma palavra por meio de alguma(s) de sua(s) sílaba(s) ou
%% letra(s)), siglas (conjunto de letras iniciais dos vocábulos e/ou números que representa um determinado nome) e acrônimos
%% (conjunto de letras iniciais dos vocábulos e/ou números que representa um determinado nome, formando uma palavra pronunciável).
%%
%% Este arquivo para definição de abreviaturas, siglas e acrônimos é utilizado com a opção "glossaries" (pacote)

%% Abreviaturas: \abreviatura{rótulo}{representação}{definição}

\abreviatura{art.}{art.}{Artigo}
\abreviatura{cap.}{cap.}{Capítulo}
\abreviatura{sec.}{sec.}{Seção}

%% Siglas: \sigla{rótulo}{representação}{definição}

\sigla{ml}{ML}{Aprendizado de Máquina, do inglês \textit{Machine Learning}}
\sigla{ia}{IA}{Inteligência Artificial}
\sigla{lgpd}{LGPD}{Lei Geral de Proteção de Dados}
\sigla{isp}{ISP}{Provedor de Serviço de Internet, do inglês \textit{Internet Service Provider}}
\sigla{json}{JSON}{Notação de Objeto JavaScript, do inglês \textit{JavaScript Object Notation}}
\sigla{rest}{REST}{Transferência de Estado Representacional, do inglês \textit{Representational State Transfer}}
\sigla{http}{HTTP}{Protocolo de Transferência de Hipertexto, do inglês \textit{HyperText Transfer Protocol}}
\sigla{csv}{CSV}{Valores Separados por Vírgula, do inglês \textit{Comma-Separated Values}}
\sigla{fes}{FES}{Estimulação Elétrica Funcional, do inglês \textit{Functional Electrical Stimulation}}
\sigla{api}{API}{Interface de Programação de Aplicações, do inglês \textit{Application Programming Interface}}
\sigla{ttl}{TTL}{Tempo de Vida, do inglês \textit{Time to Live}}
\sigla{uuid}{UUID}{Identificador Único Universal, do inglês \textit{Universally Unique Identifier}}
\sigla{jwt}{JWT}{Token Web JSON, do inglês \textit{JSON Web Token}}
\sigla{url}{URL}{Localizador Uniforme de Recursos, do inglês \textit{Uniform Resource Locator}}
\sigla{nist}{NIST}{Instituto Nacional de Padrões e Tecnologia, do inglês \textit{National Institute of Standards and Technology}}
\sigla{ux}{UX}{Experiência do Usuário, do inglês \textit{User Experience}}
\sigla{spp}{SPP}{Perfil de Porta Serial, do inglês \textit{Serial Port Profile}}
\sigla{ble}{BLE}{Bluetooth de Baixa Energia, do inglês \textit{Bluetooth Low Energy}}
\sigla{orm}{ORM}{Mapeamento Objeto-Relacional, do inglês \textit{Object-Relational Mapping}}
\sigla{pcb}{PCB}{Placa de Circuito Impresso, do inglês \textit{Printed Circuit Board}}
\sigla{ssh}{SSH}{Shell Seguro, do inglês \textit{Secure Shell}}
\sigla{sdk}{SDK}{Kit de Desenvolvimento de Software, do inglês \textit{Software Development Kit}}


%%%% GLOSSÁRIO


%% LEIA:

%% Para usar o \gls, você deve colocar a sigla aqui em \sigla

%% ADICIONAR SIGLAS: Quando você inclui alguma sigla, pode ser necessário compilar umas duas vezes para essa aparecer na lista de siglas.

%% ATENÇÃO REMOVER SIGLAS: se você remover a sigla do seu texto (não for usar mais), você deve comentar essa aqui e remover os \gls{} dessa sigla (se não vai ficar aparecendo a sigla na lista). Em caso de ERRO, quando o LaTeX informa que você ainda tem a sigla no texto, mesmo que não tenha. Você deve limpar o cache - no OverLeaf, clique no erro, vá:
%  ->view error
%%   ->(role para baixo, até o final)
%%     ->e clique em Clear cached files


%% Acrônimos: \acronimo{rótulo}{representação}{definição}
%\acronimo{gimp}{Gimp}{Programa de Manipulação de Imagem GNU, do inglês \textit{GNU Image Manipulation Program}}
